\documentclass[10pt,twocolumn]{article}
\usepackage[T1]{fontenc}
\usepackage[utf8]{inputenc}
\usepackage{lmodern}
\usepackage[margin=1.8cm,top=2.4cm,bottom=2cm,columnsep=0.6cm]{geometry}
\usepackage{fancyhdr}
\usepackage{titlesec}
\usepackage{booktabs}
\usepackage{amsmath,amssymb,amsfonts}
\usepackage{microtype}
\usepackage{xcolor}
\usepackage{mdframed}
\usepackage{parskip}
\usepackage{enumitem}
\usepackage{hyperref}

\definecolor{rulered}{RGB}{180,30,30}
\definecolor{boxgray}{RGB}{245,245,245}
\definecolor{keywordgray}{RGB}{100,100,100}

\pagestyle{fancy}
\fancyhf{}
\fancyhead[L]{\textsc{\small Claude Mag}}
\fancyhead[C]{\small\textit{Machine Learning Research Digest}}
\fancyhead[R]{\small 2026-10-14}
\fancyfoot[C]{\small\thepage}
\renewcommand{\headrulewidth}{0.8pt}

\titleformat{\section}{\normalsize\bfseries\color{rulered}}{}{0pt}{}%
  [\vspace{-4pt}\color{rulered}\rule{\columnwidth}{0.4pt}]
\titlespacing{\section}{0pt}{8pt}{4pt}

\hypersetup{colorlinks=true,urlcolor=rulered,linkcolor=black}

\begin{document}
\setlength{\parindent}{0pt}
\setlength{\parskip}{4pt}

{\small\color{keywordgray}\textbf{Keywords:}
  mechanistic interpretability \textbullet{}
  attention heads \textbullet{}
  activation patching \textbullet{}
  task-dependent circuits}\par\vspace{4pt}

{\LARGE\bfseries\raggedright
  Task Context Shapes Component Roles in Language Models}\par\vspace{4pt}

{\large\itshape\raggedright
  The Same Attention Head Does Different Jobs Depending on the Task---
  and Response Maps Predict This Redistribution Across Five Model Families}\par\vspace{6pt}

{\small\textbf{By} Sovesh Mohapatra \& Dani S Bassett}\par\vspace{2pt}

{\small\color{keywordgray}arXiv:2610.13895 \textbullet{} October 2026}
\par\vspace{2pt}\noindent{\color{rulered}\rule{\columnwidth}{0.6pt}}\vspace{6pt}

The same attention head can carry out fundamentally different computations
depending on the task context, even when the factual content of the prompt
is held constant. Across five model families (1.5--14B), switching tasks
redistributes influence among heads, and a linear \emph{response map} fitted
per task predicts a downstream head's change from an upstream activation edit,
cutting mean squared error by 53.6\% in Gemma-2-9B.

\begin{mdframed}[backgroundcolor=boxgray,linecolor=rulered,linewidth=1pt,
    innerleftmargin=6pt,innerrightmargin=6pt,
    innertopmargin=6pt,innerbottommargin=6pt]
\textbf{\small AT A GLANCE}\par\vspace{4pt}
\begin{description}[leftmargin=0pt,labelsep=4pt,itemsep=2pt,parsep=0pt,
    font=\small\bfseries]
  \item[Problem:] Mechanistic interpretability studies heads in a single-task
    context; it is unknown whether those functional roles transfer.
  \item[Why it matters:] Safety interventions derived from single-task circuits
    may fail or backfire under different task contexts.
  \item[Method:] Graph-fact prompts isolate task context from factual content;
    activation patching measures cross-task influence redistribution.
  \item[Result:] Mean cosine similarity between influence profiles across tasks
    is 0.41; response maps transfer with 50--70\% of MSE reduction preserved.
  \item[Evidence:} Five models, five task types, causal intervention validation.
\end{description}
\end{mdframed}

\section{The Big Picture}

Mechanistic interpretability has made progress identifying attention heads
responsible for specific behaviors---induction heads for copying, name-mover
heads for indirect object identification. But almost all such analyses are
conducted under a \emph{single} task context.

Mohapatra and Bassett ask: does the same head remain responsible for the same
computation when the task changes, even if the factual content of the prompt
is unchanged? Using carefully controlled \textbf{graph-fact} prompts that hold
syntactic structure and factual entities constant while varying only the task
instruction, they find that component roles are substantially reshuffled.

This challenges the view of neural circuits as fixed functional units and
suggests that safety interventions must account for task-dependent rerouting.

\section{Why Existing Approaches Fall Short}

\textbf{Single-task circuit tracing} (e.g., path patching, causal tracing)
identifies the heads responsible for a behavior in one setting but provides
no prediction of how those heads behave under a different task with identical
factual content.

\textbf{Universal probes} train on a fixed dataset to locate where facts are
stored. They miss the dynamic routing that occurs at task-switch time.

\textbf{Attention pattern visualization} shows where a head attends but not
how its output influences downstream layers. Influence requires causal
intervention, not observation.

\section{Core Mechanism: Graph Facts and Response Maps}

\textbf{Graph-fact pairs} are minimal-pair prompts that share a factual entity
(e.g., a country capital) at the same token position across tasks but differ
only in the task instruction (factual recall, translation, analogy, completion,
QA). This controls for syntactic and positional confounders.

\textbf{Influence measurement.} For each head $(h,\ell)$, the indirect effect
under task $t$ is:
\[
\mathrm{IE}^{(t)}(h,\ell) = \frac{1}{N}\sum_{i=1}^N
\!\left[p_\theta(a^{*(t)}_i|x^{(t)}_i,\mathrm{patch}_{h,\ell})
- p_\theta(a^{*(t)}_i|x^{(t)}_{\mathrm{corrupt},i})\right]
\]

\textbf{Response maps.} A linear map $R^{(t)}_{(h,\ell)\to(h',\ell')}$ predicts
the output change at downstream head $(h',\ell')$ from an activation edit at
upstream head $(h,\ell)$ under task $t$:
\[
\hat{\Delta}a_{h',\ell'} = R^{(t)}_{(h,\ell)\to(h',\ell')} \cdot \Delta a_{h,\ell}
\]

Quality is measured by fractional MSE reduction relative to a zero predictor.

\section{Technical Formulation}

Let $\mathrm{IE}^{(t)} \in \mathbb{R}^{N_{\text{heads}}}$ be the influence
vector across all heads under task $t$. The cross-task similarity:
\[
\rho(t, t') = \frac{\mathrm{IE}^{(t)} \cdot \mathrm{IE}^{(t')}}
                   {\|\mathrm{IE}^{(t)}\|\,\|\mathrm{IE}^{(t')}\|}
\]

A value of $\rho < 1$ indicates that the relative importance of heads
changes across tasks. The authors observe $\rho \in [0.18, 0.67]$ across
task pairs (mean 0.41).

Response map quality:
\[
\mathrm{MSE\text{-}red} = 1 - \frac{\|\Delta a_{h',\ell'} - \hat{\Delta}a_{h',\ell'}\|^2}
                                    {\|\Delta a_{h',\ell'}\|^2}
\]
Gemma-2-9B: 53.6\%; Llama-3.1-8B: 41.2\%; Phi-3-mini: 38.8\%; Qwen3-14B: 5.0\%.

\section{Learning or Inference Procedure}

\begin{enumerate}[itemsep=2pt,parsep=0pt]
  \item Construct graph-fact prompts for five task types; ensure factual
    entities appear at the same token position.
  \item For each model and task, run activation patching to compute
    $\mathrm{IE}^{(t)}(h,\ell)$ over all heads.
  \item Fit response maps $R^{(t)}$ on 80\% of prompts; evaluate on 20\%.
  \item Perform targeted early-layer edits; measure downstream answer-token
    probability changes.
  \item Compare influence profiles across tasks via cosine similarity.
\end{enumerate}

\section{What the Guarantee Says}

The response maps are empirically validated: for one factual intervention,
the Gemma-2-9B map reduces MSE by 53.6\% vs.\ averaging the two
task-specific maps. Cross-task transfer preserves 50--70\% of this
reduction, suggesting a shared sub-structure that is partially predictable
from single-task analysis.

\section{Experimental Findings}

\begin{itemize}[itemsep=2pt,parsep=0pt]
  \item \textbf{Models:} Phi-3-mini-3.8B, Gemma-2-9B, Llama-3.1-8B, Qwen3-14B,
    Mistral-7B.
  \item \textbf{Tasks:} Factual recall, translation, analogy, sentence
    completion, open-ended QA.
  \item Mean cross-task influence cosine similarity: \textbf{0.41}
    (range 0.18--0.67).
  \item Response map MSE reduction: Gemma-2-9B \textbf{53.6\%}; Llama-3.1-8B
    41.2\%; Qwen3-14B 5.0\%.
  \item Cross-task transfer: maps from factual recall transfer best to sentence
    completion (0.62), worst to translation (0.19).
\end{itemize}

\section{Ablations and Interpretation}

\textbf{Fixed-task vs.\ cross-task maps:} Single-task maps outperform
cross-task maps by 15--20\% MSE on the training task, confirming genuine
task-specific routing.

\textbf{Layer range:} Early-layer edits (layers 3--7 of 28) are the primary
routing bottleneck; mid-layer response maps are less predictive.

\textbf{Attention vs.\ MLP:} Attention heads account for 70--85\% of
task-dependent influence redistribution; MLP contributions are less variable
across tasks.

\textbf{Safety implication:} A circuit marked ``safe'' under factual recall may
be activated differently under translation or QA context. Cross-task analysis
is recommended before deployment.

\section*{Reference}

Sovesh Mohapatra, Dani S Bassett.
\textbf{Task Context Shapes Component Roles in Language Models}.
arXiv:2610.13895, October 2026.\\
\url{https://arxiv.org/abs/2610.13895}

\end{document}
